La introducción antecede al cuerpo del trabajo, pero se escribe de último. Constituye el primer capítulo y la página número uno. A diferencia del resumen, es más extenso, pues contextualiza y describe la composición del trabajo. Suele incluirse la siguiente información:

\begin{itemize}
    \item Campo de estudio.

    \item Finalidad del trabajo.

    \item Delimitación y definición del tema.

    \item Metodología y procedimientos, sin mayor detalle, ya que tienen su capítulo aparte.

    \item Conclusiones más importantes.
    
\end{itemize}

      

      

      

      

      